\subsection{A simultaneous basis with one additional contiguous pivot block}
\label{sec:controlled-projector-basis}

The three-class rank bound can be strengthened while keeping the finite
scalar producer unchanged. The additional structure is used only for the
stage-two auxiliary sink edges. The same rational change of coordinates
preserves the needed corner-rank conditions on the other two large classes.

Put $F=\mathbb Q^h$, $V=\mathbb Q^H$, $H=h^2$, and $m=hH$.
Reorder the tensor factors so that the original third factor is $F$ and the
first two factors form $V$. List coordinates in the order
\[
 (\alpha,i),\qquad 1\le\alpha\le h,\quad 1\le i\le H,
\]
with $i$ varying fastest. Write $\Pi_t$ and $\Pi_q$ for rank-one
idempotents, and $\Pi_U$ for a rank-$h$ idempotent. The three large
projector classes then have the following descriptions:
\begin{align*}
 A_1&=I-I_F\otimes\Pi_q-\Pi_t\otimes\Pi_U,
       &\Pi_q\Pi_U=\Pi_U\Pi_q=0,\\
 A_2&=(I_F-\Pi_t)\otimes I_V,\\
 A_3&=(I_F-\Pi_t)\otimes(I_V-\Pi_q).
\end{align*}
Their nullities are $2h$, $H$, and $H+h-1$, respectively.
For the first class, $q$ spans the line associated with
$t_B\otimes t_{\pi(A)}$, $U=F\otimes\langle t_A\rangle$, and
$t=t_B$. The orthogonality in the bank matching gives the displayed
zero products and, in particular, $q\notin U$.
The second class is the stage-two sink projection. The third is the
stage-three data entrance projection. Transposing any of these matrices
preserves its displayed form, with the corresponding dual rank-one lines
and subspaces.

Choose $K\in\operatorname{GL}(V)$ and independent matrices
$G_i\in\operatorname{GL}(F)$ for $1\le i\le H$. Define
\[
 T(u\otimes e_i)=G_i u\otimes e_i,\qquad S=T(I_F\otimes K).
\]
Thus $T$ changes the $h$ first-factor coordinates independently over each
of the $H$ multiplicity coordinates. The parameters are rational matrices;
$S$ is a change of basis in the finite address-label space, not a tape
operation that must be performed on the input streams.

\begin{lemma}[Simultaneous controlled basis]
\label{lem:controlled-projector-basis}
For the finite collection of projectors in the retained $h=28$ motif,
one can choose $K,G_1,\ldots,G_H$ over $\mathbb Q$ so that:
\begin{enumerate}
\item for every $A_1$, the submatrix of $SA_1S^{-1}$ in its first
      $2h$ rows and last $2h$ columns is nonsingular;
\item for every $A_3$, the submatrix of $SA_3S^{-1}$ in its first
      $H+h-1$ rows and last $H+h-1$ columns is nonsingular;
\item for every $A_2$, the submatrix of $SA_2S^{-1}$ in its first
      $H$ rows and last $H$ columns is diagonal and nonsingular.
\end{enumerate}
\end{lemma}
\begin{proof}
First consider the third assertion. The factor $I_F\otimes K$ commutes
with $A_2$. At multiplicity coordinate $i$, the transformed operator is
$G_i(I_F-\Pi_t)G_i^{-1}$. Consequently its first-to-last block is diagonal
in $i$, with diagonal entry
\[
 [G_i(I_F-\Pi_t)G_i^{-1}]_{1,h}.
\]
For a nonzero rank-one projector and $h\ge2$, this entry is a nonzero
rational function of $G_i$: writing $\Pi_t=t\varphi$ with $\varphi(t)=1$,
it equals $-(e_1^{\mathsf T}G_it)(\varphi G_i^{-1}e_h)$, and the two
factors can both be made nonzero. Thus the required diagonal entries can
all be required to be nonzero.

For the first assertion, let $Q_1=I-A_1$. Its image is the direct sum
\[
 \operatorname{im}Q_1=(F\otimes q)\oplus(t\otimes U).
\]
Choose a basis matrix $B$ of the $h$-space $U$. After applying $K$, write
$q'=Kq$ and $B'=KB$. The nullspace vector parametrized by $u\in F$ and
$c\in\mathbb Q^h$ has multiplicity-$i$ coordinate
\[
 G_i(q'_i u+tB'_i c).
\]
The first $2h$ coordinate rows therefore evaluate
\[
 q'_i\ell_i(u)+\ell_i(t)B'_i c,
 \qquad \ell_i=e_1^{\mathsf T}G_i,\quad 1\le i\le2h.
\]
We show that their determinant is not identically zero in the parameters.
Choose $K$ so that $q'_i\ne0$ for $1\le i<h$ and the submatrix
$[q'\ B']$ in rows $h,\ldots,2h$ is nonsingular. Such a $K$ exists because
$[q\ B]$ has column rank $h+1$ and an invertible linear map can send this
ordered independent system to any other one. On the first $h-1$ rows,
choose $\ell_i$ to be a basis of the annihilator of $t$. These rows
determine $u$ modulo $\langle t\rangle$. On the remaining $h+1$ rows,
choose one fixed functional $\ell$ with $\ell(t)=1$. After the first
$h-1$ coordinates of $u$ are determined, these rows form exactly the
nonsingular matrix $[q'\ B']$ on the remaining coordinate of $u$ and $c$.
Every prescribed nonzero $\ell_i$ extends to the first row of an invertible
$G_i$. This proves that restriction of $S\operatorname{im}Q_1$ to the
first $2h$ coordinates can have full rank.

For the second assertion, the image of $Q_3=I-A_3$ is
\[
 \operatorname{im}Q_3=(\langle t\rangle\otimes V)
          \oplus(W_t\otimes\langle q\rangle),
 \qquad W_t=\operatorname{im}(I_F-\Pi_t),\quad\dim W_t=h-1.
\]
Since the first summand contains $t\otimes V$, its image under $I\otimes K$
is unchanged. Write a vector after that change as
$t v_i+q'_i u$ at multiplicity coordinate $i$, where $v\in\mathbb Q^H$
and $u\in W_t$. The first $H$ coordinate rows give
\[
 \ell_i(t)v_i+q'_i\ell_i(u),\qquad 1\le i\le H.
\]
Choose all $\ell_i(t)\ne0$. They determine the $H$ independent coordinates
$v_i$ once $u$ is known. The next $h-1$ rows are the second rows of
$G_i$ at $i=1,\ldots,h-1$. If $f_i=e_2^{\mathsf T}G_i$, eliminating
$v_i$ gives
\[
 q'_i\left(f_i-\frac{f_i(t)}{\ell_i(t)}\ell_i\right)(u).
\]
Choose $K$ with these $q'_i$ nonzero and choose the displayed $h-1$
functionals to be a basis of the annihilator of $t$, restricted to $W_t$.
This is possible because $F=\langle t\rangle\oplus W_t$: take each
$\ell_i(t)=1$, take the relevant $f_i$ to annihilate $t$, and complete
these independent first two rows to $G_i$. Thus restriction to the
first $H+h-1$ coordinates can also have full rank.

Apply the same two arguments to the transposed projectors at the last
coordinate rows. This is a condition within the same parameter family,
because
\[
 S^{-\mathsf T}=T^{-\mathsf T}(I_F\otimes K^{-\mathsf T}),
\]
and $T^{-\mathsf T}$ acts independently by $G_i^{-\mathsf T}$.
They show that the corresponding last-coordinate restrictions of the
transformed dual nullspaces can each have full rank.

To combine the requirements, express every rank condition as a nonzero
minor after clearing the denominators of the parameter matrices and their
inverses. Each is a polynomial that is not identically zero, by the
constructions just given. Their finite product, together with the
invertibility conditions on $K$ and all $G_i$, is still a nonzero polynomial.
A nonzero rational polynomial cannot vanish at every rational tuple, so
one rational choice meets all conditions simultaneously. The individual
witness choices above establish that the polynomials are nonzero; they
need not themselves meet all the other conditions.

Finally, for a rank-$r$ nullspace projector $Q=UV$ with $VU=I_r$, its
first-$r$ by last-$r$ submatrix after conjugation factors as the product
of the two square restriction matrices just considered. Both are
nonsingular. The corresponding block of $I-SQS^{-1}$ is its negative,
since these first and last index sets are disjoint. This proves the first
two assertions. The third has already been proved directly.
\end{proof}

The argument is constructive in the finite-data sense required by the
network interface: enumerate rational parameter matrices and retain the
first choice passing the finitely many specified determinant tests. The
lemma proves termination. All parameters and tests depend only on the fixed
motif, not on the eventual multiplication input size. Conjugate every
terminal and gate matrix by this one $S$. Their endpoint differences remain
$I_m$, common gate matrices remain common, and all edge ranks are preserved.

\paragraph{The improved stage-two pivot profile.}
The lower-triangular elimination for a large projector first eliminates its
nonsingular corner block. For $A_2$, this block is already diagonal and its
$H$ pivots lie in matching physical order in two contiguous sets of fields.
After fixed diagonal normalizations, they are one $H\times H$ identity
block. The Schur-complement argument in the large-projector lemma leaves
the middle identity block of size $m-2H$. Hence a stage-two sink edge is
implemented by two contiguous child interchanges, of widths
\[
 Hb\quad\hbox{and}\quad(m-2H)b,
\]
with no singleton recursive calls on its corner. Its rank remains
$H+(m-2H)=m-H$, so the rank sum at exponent one is unchanged.
The first and third edge classes retain their previous singleton-corner
and contiguous-middle implementations by the first two parts of the lemma.

\paragraph{Exact strengthened moment.}
Let $B=v^2R$ and let $S_0$ be the singleton count from the three-class
certificate. Replace $S_0$ by $S_0-BH$ and add $B$ children of width $H/m$.
For $h=28$ the new singleton count is
\[
 S_1=105555640096680883200.
\]
The normalized child widths and their rank-mass weights are respectively
\[
 \left(\frac1{21952},\frac{195}{196},\frac{13}{14},
             \frac{10165}{10976},\frac1{28}\right)
\]
and
\[
 \left(\frac{10841531}{520019360},\frac{12827685}{26000968},
       \frac{855179}{1857212},\frac{20865}{2736944},
       \frac{65783}{3714424}\right).
\]
Use the four preceding rational logarithm bounds and additionally
$\log 28<33323/10000$. With
\[
 a_{\rm b}=\frac{246}{10^9},\qquad\tau=1-a_{\rm b},
\]
the same bound $e^x\le(1-x)^{-1}$ gives a strict rank-moment gap exceeding
$4.76\times10^{-11}$. The exact rational calculation is in
\texttt{scripts/controlled\_bit\_rank\_moment.py}. This proves the stronger
interchange saving using the unchanged PR~\#7 scalar producer and the
same arbitrary-width recurrence interface.
